\documentclass[11pt]{article}
\usepackage[a4paper,margin=2cm]{geometry}
\usepackage{amsmath,amssymb,graphicx,booktabs}
\begin{document}
\section*{Validation of the Mixed-Zonal GEqOE Short-Period Map}

This note validates the first-order mixed-zonal GEqOE short-period map for the
truncated zonal set $\{J_2,J_3,J_4,J_5\}$. The mean state is propagated in
$(\nu,p_1,p_2,M,q_1,q_2)$, where $M=L-\Psi$ is the uniformly advancing mean
fast phase. Osculating GEqOE reconstruction is obtained by applying the exact
degree-wise short-period corrections in $(g,Q,\Psi,\Omega,M)$ and then solving
the generalized Kepler equation for $K$.

\subsection*{Osculating-history reconstruction}
\begin{center}
\begin{tabular}{lcccccc}
\toprule
case & $M_0$ [deg] & $K$ RMS [rad] & pos RMS [km] & pos max [km] & $p_1$ rel RMS & $q_1$ rel RMS \\
\midrule
low-e & 20 & 1.026e-05 & 1.070e-01 & 1.861e-01 & 3.225e-06 & 1.240e-05 \\
low-e & 140 & 1.030e-05 & 1.070e-01 & 1.935e-01 & 4.171e-06 & 1.194e-05 \\
high-e & 35 & 2.408e-06 & 4.220e-02 & 1.088e-01 & 5.339e-07 & 4.653e-06 \\
high-e & 170 & 2.369e-06 & 4.797e-02 & 1.229e-01 & 4.756e-07 & 4.985e-06 \\
\bottomrule
\end{tabular}
\end{center}

The representative component histories are shown in
Figure~\ref{fig:low} for the low-eccentricity case and in
Figure~\ref{fig:high} for the high-eccentricity case. Figure~\ref{fig:cart}
collects the Cartesian position errors for all anomaly cases.

\begin{figure}[p]
\centering
\includegraphics[width=0.95\textwidth]{docs/geqoe_averaged/figures/zonal_short_period_lowe_components.png}
\caption{Low-eccentricity osculating GEqOE and Cartesian reconstruction.}
\label{fig:low}
\end{figure}

\begin{figure}[p]
\centering
\includegraphics[width=0.95\textwidth]{docs/geqoe_averaged/figures/zonal_short_period_highe_components.png}
\caption{High-eccentricity osculating GEqOE and Cartesian reconstruction.}
\label{fig:high}
\end{figure}

\begin{figure}[p]
\centering
\includegraphics[width=0.90\textwidth]{docs/geqoe_averaged/figures/zonal_short_period_cartesian_errors.png}
\caption{Cartesian position errors across the anomaly sweep.}
\label{fig:cart}
\end{figure}

\subsection*{Zonal scaling}
For the representative high-eccentricity case, the position RMS reconstruction
error was recomputed under zonal scale factors
$\{1.00, 0.50, 0.25\}$.
The fitted log-log slope was
\[
\frac{d\log(\mathrm{RMS})}{d\log(\lambda)} = 2.000,
\]
which is consistent with the expected second-order residual of a first-order
mean-plus-short-period model.

\begin{figure}[h!]
\centering
\includegraphics[width=0.70\textwidth]{docs/geqoe_averaged/figures/zonal_short_period_scaling.png}
\caption{RMS Cartesian reconstruction error versus zonal scale factor.}
\end{figure}

\subsection*{Conclusion}
Within the truncated mixed-zonal problem, the GEqOE averaged formulation is now
closed to first order: the exact symbolic mean drift and the exact symbolic
short-period map together reconstruct the full osculating GEqOE and Cartesian
state with small residuals, and those residuals scale as expected for a
first-order theory.

\end{document}
